% Nouveau EGA V (cote 25): archive cote 25, site page 144 (PDF page 145).
% English translation of the checked French annotated master.
% AI-translated by GPT-6 Astra (OpenAI, ChatGPT Pro mode).
\scanpage{25/144}
\mspage{\hand{5}}

The column comes from the transitivity exact sequence for the
smooth morphisms $X\hand{\to} Y$ and $Y\hand{\to} S$, by \var{\xout{xx}}{}{typewriter strike-out}{ega25_25_144-v0001} tensoring with $\underline{O}_{V1}$ (this remains
exact, since all the modules in the sequence are loc free). The horizontal
row is a particular case of an exact sequence obtained whenever,
on X \var{\deleted{\var{\xout{smooth}}{}{typewriter strike-out}{ega25_25_144-v0004}}}{}{deletion}{ega25_25_144-v0003} over S, one has a section \hand{$\mu$} of a loc free module F
and takes the scheme W of zeros; one finds an exact sequence
\[ \overset{\hand{\vee}}{F}\otimes \underline{O}_X \hand{\longrightarrow} \underline{\hand{\Omega}}^1_{X/S}\otimes \underline{O}_W \hand{\longrightarrow} \underline{\hand{\Omega}}^1_{W/S} \hand{\longrightarrow} 0 \]
and \var{\inserted[typed]{(if X/S is smooth)}}{(if X/S is smooth)}{insertion}{ega25_25_144-v0005} the first homomorphism is injective exactly at the points where
W is smooth over X with the ``right'' relative \unsure{dimension} (i.e. everywhere in
\var{\inserted[typed]{\var{\xout{practically}}{}{typewriter strike-out}{ega25_25_144-v0007}}}{\var{\xout{practically}}{}{typewriter strike-out}{ega25_25_144-v0007}}{insertion}{ega25_25_144-v0006}
the present case). \var{\xout{This homomorphism}}{}{typewriter strike-out}{ega25_25_144-v0008} This exact sequence is \var{\xout{written in}}{}{typewriter strike-out}{ega25_25_144-v0009}
\var{\xout{section 16 in the case where F = $\underline{O}_X$, the generalization}}{}{typewriter strike-out}{ega25_25_144-v0010} \var{\xout{(except that}}{}{typewriter strike-out}{ega25_25_144-v0011} an
immediate consequence of the exact sequence
\[ \underline{J}/\underline{J}^2\hand{\longrightarrow} \underline{\hand{\Omega}}^1_{X/S}\otimes \underline{O}_W \hand{\longrightarrow} \underline{\hand{\Omega}}^1_{X/S} \hand{\longrightarrow} 0 \]
\var{\margin{16 \textbar}}{\placedmargin{16 \textbar}}{marginal note placed}{ega25_25_144-v0012} that appears in section 16 (one could put the version
mentioned here there as a corollary). The characterization of the set of points where
one can put a zero on the left is contained in the Jacobian criterion.

Observe that we have a canonical isomorphism
\[ \underline{P}^1_{X/S} = \underline{\hand{\Omega}}^1_{X/S} + \underline{O}_X \ ,\ \text{d'où}\ \overset{\hand{\scriptscriptstyle\vee}}{\underline{P}}^1_{X/S} = \hand{\mathcal{T}}_{X/S} + \underline{O}_X \ , \]
and on the other hand one \unsure{verifies} that the \var{\inserted[typed]{composite}}{composite}{insertion}{ega25_25_144-v0014} homomorphism \hand{$\mu$} of diagram \hand{\unsure{(D)|(1)}} is zero
on the factor $\underline{O}_{V1}$, and on the factor $\hand{\mathcal{T}}_{X/S}\otimes\underline{O}_{V1}$ \var{\inserted{it}}{it}{insertion}{ega25_25_144-v0015} reduces to
the homomorphism \var{\inserted[typed]{$M(\hand{\phi})'$}}{$M(\hand{\phi})'$}{insertion}{ega25_25_144-v0016} obtained from the section $M(\hand{\phi})$ of $\underline{\mathrm{Sym}}^2(\hand{\underline{\hand{\Omega}}}^1_{X/S})\otimes \underline{O}_{V1}$ already
considered. Thus at a point x of X, $M(\hand{\phi})$ is nondegenerate, i.e. $M(\hand{\phi})'$
is surjective iff \hand{$\mu$} is surjective \var{\inserted[typed]{at x}}{at x}{insertion}{ega25_25_144-v0017}, and we see from diagram \hand{(D)}
that this is also equivalent to saying that \hand{$\nu$} is surjective \var{\inserted[typed]{at x}}{at x}{insertion}{ega25_25_144-v0018} (since both
mean that the canonical homomorphism from\var{\xout{xx two}}{}{typewriter strike-out}{ega25_25_144-v0020} the sum of the two sub-Modules
considered in $\underline{\hand{\Omega}}^1_{X/S}\otimes \underline{O}_{V1}$ into the latter is surjective at
x). We thus find:

\mspage{\hand{144}}
\tnote{25-144-1}{Loc abbreviates locally. The source spellings transtitivité and exacten are translated without invented English spelling errors. The canceled e in une, the j-to-J capitalization of jacobien, and the added plural ending nt in signifient have no English effect and are recorded as French-only variants. The live doubt vérifie/vérifié remains on verifies. In the literal display, d'où means whence. The source's W smooth over X and the repeated X/S index in the second exact sequence are retained.}
